% Part: incompleteness
% Chapter: arithmetization-syntax
% Section: substitution

\documentclass[../../../include/open-logic-section]{subfiles}

\begin{document}

\olfileid{inc}{art}{sub}
\olsection{પ્રતિસ્થાપન}

યાદ કરો કે પ્રતિસ્થાપન એ
!!a{formula}~$!A$માં !!a{variable}~$u$ની
બધી મુક્ત ઘટનાઓની જગ્યાએ પદ~$t$
મૂકવાની ક્રિયા છે; તેને~$\Subst{!A}{t}{u}$
લખીએ છીએ. !!{variable}s, !!{formula}s
અને પદોની ગડેલ-સંખ્યાઓ પર કરવામાં આવે
ત્યારે આ ક્રિયા આદિમ પુનરાવર્તી છે.

\begin{prop}\ollabel{prop:subst-primrec}
એવું આદિમ પુનરાવર્તી વિધેય
$\fn{Subst}(x, y, z)$ છે કે
\[
\fn{Subst}(\Gn{!A}, \Gn{t}, \Gn{u}) = \Gn{\Subst{!A}{t}{u}}.
\]
\end{prop}

\begin{proof}
આદિમ પુનરાવર્તન વડે વિધેય~$\fn{hSubst}$
આ પ્રમાણે વ્યાખ્યાયિત કરી શકીએ:
\begin{multline*}
\begin{aligned}
\fn{hSubst}(x, y, z, 0) & = \emptyseq \\
\fn{hSubst}(x, y, z, i+1) & =
\end{aligned}\\
\begin{cases}
\fn{hSubst}(x, y, z, i) \concat y & \text{જો $\fn{FreeOcc}(x, z, i)$ હોય} \\
\fn{append}(\fn{hSubst}(x, y, z, i), (x)_{i}) & \text{અન્યથા.}
\end{cases}
\end{multline*}
હવે~$\fn{Subst}(x, y, z)$ને
$\fn{hSubst}(x, y, z, \len{x})$ તરીકે
વ્યાખ્યાયિત કરી શકીએ.
\end{proof}

\begin{prop}
\ollabel{prop:free-for}
ગડેલ-સંખ્યા~$x$ ધરાવતા સૂત્રમાં
ગડેલ-સંખ્યા~$y$ ધરાવતું પદ,
ગડેલ-સંખ્યા~$z$ ધરાવતા ચલની જગ્યાએ
!!{free for} હોય ત્યારે અને માત્ર ત્યારે
સત્ય થતો સંબંધ~$\fn{FreeFor}(x, y, z)$
આદિમ પુનરાવર્તી છે.
\end{prop}

\begin{proof} અભ્યાસપ્રશ્ન. \end{proof}

\begin{prob}
\olref[inc][art][sub]{prop:free-for} સાબિત કરો.
\end{prob}

\end{document}
